\subsubsection{Arbitrary fixed boundary on a surface mapping to a curve}

\begin{proposition}[Proved: the logarithmic fibration criterion]
\label{D:log-fibration-criterion}
Let $k$ be algebraically closed, let $S$ be a normal connected
proper surface, and let $f:S\to B$ be a morphism to a smooth
connected proper curve.  Let $Y\subsetneq S$ be closed and put
$U=S\setminus Y$.  Suppose that there is a nonempty open
$B^\circ\subset B$ such that $S^\circ=f^{-1}(B^\circ)\to B^\circ$
is a smooth proper family of geometrically connected curves,
and $(Y\cap S^\circ)_{\mathrm{red}}$ is a finite \'etale relative
divisor $D^\circ$ in that family, possibly empty.  Then the
fixed-boundary tame group $Q(S,Y)$ is topologically finitely
generated.  The boundary may contain whole fibres, and there
is no restriction on the multiplicities of the bad fibres.

More precisely, let $N$ be the image in $Q(S,Y)$ of the tame
fundamental group of the geometric generic pointed fibre.
Then $N$ is closed, normal, and finitely generated, and
$Q(S,Y)/N$ is a quotient of $\pi_1(B^\circ)$.  In characteristic
$p>0$, the image in this quotient of the wild inertia at any
$b\in B\setminus B^\circ$ is finite, of order dividing the
$p$-part of the gcd of the multiplicities of all components
of $S_b$.
\end{proposition}
\begin{proof}
Give $B^\circ$ the trivial log structure and give $S^\circ$
the log structure of the relative divisor $D^\circ$.
This is a proper log smooth morphism of fs log schemes, its
base is log regular, and its underlying fibres are geometrically
connected.  The logarithmic homotopy theorem gives an exact
sequence
\[
 \pi_1^{\mathrm t}\bigl(S_{\overline\eta}\setminus
                       D_{\overline\eta}\bigr)
 \longrightarrow \pi_1(S^{\circ,\log})
 \longrightarrow \pi_1(B^\circ)\longrightarrow1.
\]
Here the base geometric point has trivial log structure, so
the reduced-covering-point limit in the theorem has the
ordinary geometric point as a terminal object.  We use
\cite[Proposition 2.1 and Theorem 2.3, pp. 13--16]{Hoshi2006B}.
The middle group classifies finite \'etale covers of
$U^\circ=S^\circ\setminus D^\circ$ tame at the codimension-one
points of $D^\circ$, by
\cite[Proposition A.10, pp. 57--58]{Hoshi2006B}.
Thus its notion of tameness is exactly the one imposed by
the fixed-boundary quotient.

Restriction of a connected finite \'etale cover of the normal
integral $U$ to the dense open $U^\circ$ remains connected.
Every cover counted by $Q(S,Y)$ is tame along $D^\circ$.
Consequently the natural homomorphism
\[
             \pi_1(S^{\circ,\log})\longrightarrow G:=Q(S,Y)
\]
is surjective.  The image $N$ of the first group in the exact
sequence is therefore closed and normal in $G$, and there is
a surjection
\[
                    \pi_1(B^\circ)\longrightarrow Q:=G/N.
\]
The tame fundamental group of a smooth geometric curve with
finitely many punctures is finitely generated
\cite[XIII, Corollary 2.12, pp. 392--393]{SGA1}.
Hence $N$ is finitely generated.

We prove that $Q$ has finite wild inertia at every omitted
base point.  Let $Q_0$ be any finite quotient of $Q$.
Its connected torsor $V\to U$ is tame along the original
boundary $Y$, and its restriction to $U^\circ$ comes from a
connected finite Galois cover $B'^\circ\to B^\circ$ with group
$Q_0$.  Normalize $B$ in its function field to obtain the
finite map $B'\to B$.

Fix $b\in B\setminus B^\circ$ and a component $E$ of the
fibre $S_b$ with multiplicity $m_E$.  Let $w$ be any
prolongation of its discrete valuation to $k(V)$, and let
$e_E$ be the ramification index over the $E$-valuation of
$k(S)$.  If the generic point of $E$ belongs to $U$, then
$e_E=1$, since $V\to U$ is \'etale.  Otherwise $E$ is an
original boundary divisor, so $e_E$ is prime to $p$ in positive
characteristic, since this cover is tame there.  Restrict $w$
to the intermediate field $k(B')$.  Multiplication of indices
of the discrete value groups in the two towers gives
\[
                  e(b'/b)\mid m_E e_E.
\]
In characteristic $p>0$, this implies
\[
                v_p\bigl(e(b'/b)\bigr)\leq v_p(m_E)
                \qquad\hbox{for every component }E.
\]
All closed-point residue fields of $B'$ equal $k$.  For this
finite Galois cover of smooth curves the inertia order is
therefore $e(b'/b)$, and the image of wild inertia has order
equal to its $p$-part.  Thus the image of the wild inertia
group $P_b$ in every finite quotient of $Q$ has order dividing
the $p$-part of $\gcd_E m_E$.  Its image $\overline P_b$ in
$Q$ is finite with this same bound: a profinite group whose
finite quotient orders are bounded is finite, as follows by
choosing a finite quotient of maximal order and comparing
all refinements.  In characteristic zero the groups $P_b$
are trivial, and no bound is needed.

There are only finitely many omitted base points.  Choose an
open normal subgroup $Q'\subset Q$ meeting all the finite
groups $\overline P_b$ trivially.  Let
$B_1^\circ\to B^\circ$ be the connected finite \'etale cover
corresponding to the inverse image of $Q'$ in $\pi_1(B^\circ)$,
and let $B_1$ be its smooth proper compactification.
At a point of $B_1\setminus B_1^\circ$, its wild inertia maps
into a conjugate of $P_b$ in $\pi_1(B^\circ)$: its image is
pro-$p$, whereas the quotient of the original inertia by
$P_b$ is pro-prime-to-$p$.  Its further image in $Q'$ is
therefore trivial.  Thus the surjection
\[
                    \pi_1(B_1^\circ)\longrightarrow Q'
\]
factors through $\pi_1^{\mathrm t}(B_1^\circ)$, which is
finitely generated by the same tame-curve theorem.
Consequently $Q'$ and its finite extension $Q$ are finitely
generated.  Since $N$ is finitely generated, so is $G$.
\end{proof}

\begin{theorem}[Proved: every fixed boundary on a proper surface mapping to a curve]
\label{D:surface-map-all-boundaries}
Let $S$ be a normal connected proper surface over an
algebraically closed field.  If $S$ admits a nonconstant
morphism to a proper curve, then $Q(S,Y)$ is topologically
finitely generated for every proper closed subset $Y\subset S$.
In particular, arbitrary horizontal and vertical original
boundary divisors are allowed, with no separability assumption
on the horizontal divisors and no smoothness assumption on
the original generic fibre.
\end{theorem}
\begin{proof}
Take Stein factorization of the nonconstant morphism to its
image.  As in Theorem~\ref{D:surface-map-to-curve}, we obtain
$f:S\to B$ with $B$ a smooth connected proper curve,
$f_*\mathcal O_S=\mathcal O_B$, and $F=k(B)$ algebraically
closed in $K=k(S)$.  The p-rank-one argument in that theorem
shows that $K/F$ is regular.  Hence the normal proper generic
curve $C=S_\eta$ is geometrically integral, although it need
not be smooth.

In positive characteristic, Lemma 33.27.4, tag 0BXT, gives
a finite purely inseparable extension $F_1/F$ for which the
normalization of $C_{F_1}$ is geometrically normal
\cite{C_StacksLogFibration}.  A geometrically normal curve
is smooth: after algebraic closure it is a normal curve,
hence regular and smooth by Lemma 33.43.8, tag 0B8Y, and
smoothness descends under field extensions.  Let $B_1$ be the
normalization of $B$ in $F_1$, and let $S_1$ be the normalization
of $S$ in $KF_1$.  Normalizations here are finite, since all
the schemes are of finite type over a field.  The induced
morphism $S_1\to B_1$ is proper, and its generic curve is
the smooth geometrically integral normalization just chosen.
In characteristic zero we take $F_1=F$, $B_1=B$, and $S_1=S$;
the normal generic curve is already smooth.

Let $E_1,\ldots,E_s$ be the residue fields of the finitely many
closed points cut out by the reduced horizontal boundary
on the generic curve of $S_1/B_1$.  In positive characteristic
choose $r$ sufficiently large that
\[
        E_iF_2/F_2\text{ is separable for every }i,
        \qquad F_2=F_1^{1/p^r}.
\]
This is a finite purely inseparable extension, since $F_1$
has transcendence degree one over the perfect field $k$.
To verify the assertion about the $E_i$, let $E_i^{\mathrm{sep}}$
be the maximal separable subextension of $E_i/F_1$ and choose
$r$ with $E_i^{p^r}\subset E_i^{\mathrm{sep}}$ for all $i$.
Then $E_iF_2$ is a subfield of
$(E_i^{\mathrm{sep}})^{1/p^r}$, which is separable over $F_2$.
In characteristic zero no further extension is needed.

Normalize $B_1$ in $F_2$ and $S_1$ in $KF_2$, obtaining
$B_2$ and $S_2$.  Since the generic curve of $S_1/B_1$ is
already smooth, its base change to $F_2$ is normal; the second
normalization does not change that generic curve.  Its reduced
horizontal boundary therefore has exactly the separable
residue fields just described.  Shrinking a nonempty open
$B_2^\circ\subset B_2$, we may arrange that $S_2\to B_2$ is
a smooth proper geometrically connected curve family there,
that all vertical boundary lies outside it, and that the
reduced horizontal boundary is a finite \'etale relative
divisor.  For the last assertion, finite separability gives
an \'etale dense open on each horizontal boundary curve;
properness makes the horizontal map finite after shrinking,
and we remove its finitely many singularities and pairwise
intersections.  In the smooth curve family the resulting
finite \'etale closed subscheme is a relative Cartier divisor.
Thus Proposition~\ref{D:log-fibration-criterion} applies to
$(S_2,Y_2)$, where $Y_2$ is the inverse image of $Y$.

We finish by checking descent to the original fixed-boundary
quotient.  The composite $S_2\to S$ is finite and radicial.
Indeed, $KF_2/K$ is purely inseparable of some bounded
exponent $q$, and on each normal affine open $\operatorname{Spec}A$
of $S$, every element $a$ of its integral closure in $KF_2$
satisfies $a^q\in K$.  This power is integral over $A$ and
therefore belongs to $A$.  Thus the integral map is universally
injective, and lying over makes it surjective.  It is a
universal homeomorphism.  The same holds for
$U_2=S_2\setminus Y_2\to U=S\setminus Y$.
Topological invariance of the \'etale site gives
\[
                         \pi_1(U_2)\simeq\pi_1(U)
\]
\cite[Theorem 59.45.2, tag 04DY]{C_StacksLogFibration}.
Each codimension-one point of $Y_2$ lies above a
codimension-one point of $Y$, because the finite universal
homeomorphism preserves chains of specializations.  Both
local rings are DVRs.  Tame finite extensions of DVRs remain
tame after arbitrary extension of the base DVR
\cite[Lemma 15.116.9, tag 0EXY; Lemma 58.31.1, tag 0EYD]
{C_StacksLogFibration}.
Consequently every cover counted by $Q(S,Y)$ pulls back to
a cover counted by $Q(S_2,Y_2)$.  Under the displayed
identification this gives a surjection
\[
                         Q(S_2,Y_2)\longrightarrow Q(S,Y).
\]
The source is finitely generated by the proposition, so the
target is finitely generated as well.
\end{proof}

\begin{remark}
The theorem uses a morphism from the given proper surface to
a curve.  Replacing it by a rational pencil on a projective
resolution does not establish the unrestricted surface case:
an exceptional curve over an omitted point may be horizontal
for that pencil.  The original fixed-boundary condition puts
no tameness requirement on such a new horizontal valuation,
so the logarithmic homotopy sequence above does not provide
a finitely generated fibre image for that larger quotient.
\end{remark}
